\scanpage{97}
\deleted{dans la catégorie $C'(k)$ aux foncteurs compatibles avec les
limites inductives.}

On constate que le foncteur $\mathcal N$ qu'on vient de définir est exact
à gauche, \unsure{et} transforme groupes comm. en groupes comm., d'où
$C(S)\longrightarrow C'(k)$. On pose aussi
\[
\mathcal N(F)=\underline H^0_{S/k}(F)
\]
et \deleted{\unsure{de}} par suite
\[
R^i\mathcal N(F)=\underline H^i_{S/k}(F).
\]
Notons l'isomorphisme
\[
H^0_S\simeq H^0_k\circ\underline H^0_{S/k}
\]
et d'autre part que $\underline H^0_{S/k}$ \uline{transforme injectifs en injectifs}
\margin{à vérifier\ldots}
\deleted{\unsure{car on va voir} en construisant un adjoint à
$\mathcal N=\underline H^0_{S/k}$ et constatant que ce dernier est exact
à gauche :
\[
\operatorname{Hom}(\Phi,\mathcal N(F))\simeq
\operatorname{Hom}(\mathcal N^*(\Phi),F).
\]}
D'où une suite spectrale de Leray
\[
H^*(S,F)\Longleftarrow H^p(k,\underline H^q_{S/k}(F)).
\]
\mspage{96}

